% !TeX root = ../xdupgthesis_template_lc.tex

\ExplSyntaxOn
\dim_new:N \l__blind_bio_label_dim
\dim_new:N \l__blind_bio_labelsep_dim
\dim_set:Nn \l__blind_bio_label_dim { 1.4em }
\dim_set:Nn \l__blind_bio_labelsep_dim { 0.55em }
\cs_new_protected:Npn \__blind_bio_item:nn #1#2
  {
    \par
    \noindent
    \hangindent=\dim_eval:n
      { \l__xdu_bio_indent_dim + \l__blind_bio_label_dim + \l__blind_bio_labelsep_dim }
    \hangafter=1
    \raggedright
    \skip_horizontal:N \l__xdu_bio_indent_dim
    \makebox[\dim_use:N \l__blind_bio_label_dim][l]{[#1]}
    \skip_horizontal:N \l__blind_bio_labelsep_dim
    #2
    \par
  }
\cs_new_protected:Npn \bioresitem #1#2
  { \__blind_bio_item:nn {#1} {#2} }
\ExplSyntaxOff

\section{基本情况}

\section{教育背景}

\section{攻读硕士学位期间的研究成果}

% 如需保留“发表学术论文”小节，可按模板补回：
% \subsection{发表学术论文}
% \begin{bioreslist}
% \item XXX, XXX, XXX. 论文题目[J]. 期刊名, 年, 卷(期): 起止页码.
% \end{bioreslist}

\subsection{申请（授权）专利、软件著作权}
\bioresitem{1}{XXX 等．一种低功耗地磁车辆检测方法及系统：中国，申请号：2026\allowbreak1023\allowbreak90620，申请日期：2026 年 02 月 28 日（共 6 位发明人，本人为第二发明人，学生排名第一）。}
\bioresitem{2}{XXX 等．面向多次漏检的双车道外侧地磁车辆轨迹连续重构方法及系统：中国，申请号：2026\allowbreak1023\allowbreak90508，申请日期：2026 年 02 月 28 日（共 6 位发明人，本人为第三发明人，学生排名第二）。}
\bioresitem{3}{XXX 等．单地磁异常停车检测与事件上报软件 V1.0：中国，受理号：2026\allowbreak R11S0463698，申请日期：2026 年 02 月 03 日（共 5 位作者，本人排名第二，学生排名第一）。}
\bioresitem{4}{XXX 等．基于地磁传感器的拥堵车辆检测仿真系统 V1.0：中国，受理号：2026R11\allowbreak S0463768，申请日期：2026 年 02 月 03 日（共 6 位作者，本人排名第三，学生排名第二）。}
\bioresitem{5}{XXX 等．基于深度学习的小车换道意图识别和轨迹预测系统 V1.0：中国，受理号：2026R11\allowbreak S0463705，申请日期：2026 年 02 月 04 日（共 7 位作者，本人排名第四，学生排名第三）。}

\subsection{参与科研项目及获奖}
\bioresitem{1}{国家自然科学基金重点支持项目，车联网环境的智能汽车多模态高精度融合定位（U21A20446），2022.01-2025.12，已结题，负责地磁组车辆检测算法的开发与研究。}
\bioresitem{2}{荣获 2023 年度至 2024 年度西安电子科技大学优秀研究生荣誉称号。}
\bioresitem{3}{荣获 2023 年度至 2024 年度西安电子科技大学二等奖学金。}
\bioresitem{4}{荣获 2024 年度至 2025 年度西安电子科技大学二等奖学金。}
\bioresitem{5}{荣获 2025 年度至 2026 年度西安电子科技大学三等奖学金。}
